%\documentclass[11pt,a5paper]{article}
\documentclass[11pt,border=5pt,varwidth=120mm]{standalone}
\usepackage{amsmath,parskip,enumitem,pgfplots}
\pgfplotsset{compat=newest}
%\usepgfplotslibrary{patchplots}% for "patch refines"
%\usepackage{pgfplotstable} % for some table plots
% view stereo with cross eyes, or vice-versa
%\newcommand{\qview}[2]{\foreach \q in {#2,#1}}
\newcommand{\answer}[1]{\paragraph{Answer} #1} % maybe answers
\renewcommand{\vec}[1]{\text{\boldmath$#1$}}
\renewcommand{\Vec}[1]{%
  \expandafter\def\csname#1v\endcsname%
  {\ensuremath{\vec #1}}}


\title{Exercise \jobname}
\author{A. J. Roberts, \today}
\date{}

\begin{document}

\maketitle

\Vec x\Vec u\Vec v
Each plot below shows (unit) direction vectors~\vv\  (left), and the corresponding vectors~\(A\vv\) (right), for some specific matrix~\(A\). 
Estimate which \emph{pair of orthogonal} directions~\(\vv_1,\vv_2\)\ give a \emph{pair of orthogonal} result vectors~\(A\vv_1,A\vv_2\).  
By imagining this resultant pair as \(\sigma_1\uv_1,\sigma_2\uv_2\) for orthogonal unit vectors~\(\uv_1,\uv_2\), estimate the singular values~\(\sigma_1,\sigma_2\) of the matrix~\(A\).

\pgfplotsset{colormap={mySix}{
    color(0cm)=(red); 
    color(1cm)=(orange); 
    color(2cm)=(yellow!85!black); 
    color(3cm)=(green!85!black); 
    color(4cm)=(cyan); 
    color(5cm)=(blue)
    }}


\newcommand{\SVDrose}[4]{%
%  \begin{tikzpicture}%
%    \begin{axis}[footnotesize,font=\footnotesize
%        ,thick,axis equal image, axis lines=middle
%        ,colormap name=mySix ,samples=12]% 8,10,12
%        \addplot[domain=0:165.001, 
%          point meta={mod(\x,90)},
%          quiver={u=cos(\x),v=sin(\x),
%            colored=mapped color}, 
%          -stealth]    ({0},{0});
%        \addplot[domain=0:165.001, 
%          point meta={mod(\x,90)},
%          quiver={u=-cos(\x),v=-sin(\x),
%            colored=mapped color}, 
%          -stealth, densely dashed]    ({0},{0});
%        \addplot[domain=0:360,samples=73
%          ,black ,dotted] 
%          ({cos(\x)},{sin(\x)});
%    \end{axis}
%  \end{tikzpicture}
  \begin{tikzpicture}%
    \begin{axis}[small,font=\footnotesize
        ,thick,axis equal image, axis lines=middle
        ,colormap name=mySix ,samples=12% 8,10,12
        ,title={\pgfmathprintnumber[fixed]{#1}, 
                \pgfmathprintnumber[fixed]{#2}, 
                \pgfmathprintnumber[fixed]{#3}, 
                \pgfmathprintnumber[fixed]{#4} }]
        \addplot[domain=0:165.001, %=180*(1-1/samples)
          point meta={mod(\x,90)},
          quiver={u=#1*cos(\x)+#2*sin(\x), 
              v=#3*cos(\x)+#4*sin(\x),
              colored=mapped color},
         -stealth]    ({0},{0});
        \addplot[domain=0:165.001, 
          point meta={mod(\x,90)},
          quiver={u=-#1*cos(\x)-#2*sin(\x), 
              v=-#3*cos(\x)-#4*sin(\x),
              colored=mapped color},
         -stealth, densely dashed]    ({0},{0});
        \addplot[domain=0:360,samples=73
          ,black ,dotted] 
          ({#1*cos(\x)+#2*sin(\x)}
          ,{#3*cos(\x)+#4*sin(\x)});
    \end{axis}
  \end{tikzpicture}}

\edef\abcd{1}%fudged
\SVDrose1001
\SVDrose{1}{1}{0}{1}
\answer{cyan \(\vv_1 \approx\pm(0.5,0.8)\) and \(\vv_2 \approx\pm(-0.8,0.5)\), map to cyan \(A\vv_1 \approx\pm(1.4,0.9) \approx \pm(0.8,0.5)\cdot1.6\) and \(A\vv_2 \approx\pm(-0.3,0.5) \approx\pm(-0.5,0.8)\cdot0.6\).  Thus singular values of~\(A\) are~\(\sigma_1\approx1.6\) and \(\sigma_2\approx 0.6\)\,.}

\raggedright

\def\SVDroseRand{%
\pgfmathparse{round(16*rnd+4)/10}\edef\a{\pgfmathresult}
\pgfmathparse{round(10*rand)/10} \edef\b{\pgfmathresult}
\pgfmathparse{round(10*rand)/10} \edef\c{\pgfmathresult}
\pgfmathparse{round(16*rnd+4)/10}\edef\d{\pgfmathresult}
\pgfmathparse{\a*\d-\b*\c}\edef\abcd{\pgfmathresult}
\SVDrose{\a}{\b}{\c}{\d}}

\SVDroseRand
\SVDroseRand
\SVDroseRand
\SVDroseRand
\SVDroseRand
\SVDroseRand

\end{document}
